\chapter{Introducción}

\section{Objetivo del curso}
El objetivo de estas notas es registrar de forma detallada y rigurosa los conceptos, deducciones y algoritmos vistos en el curso de Mecánica Celeste (2026-1), con énfasis en la comprensión profunda de los fundamentos físicos y matemáticos. En esta clase se aborda el problema relativo de dos cuerpos, pieza clave para reducir el problema de \(N\) cuerpos a un sistema tratable.

\section{Convenciones de apuntes}
\begin{itemize}[leftmargin=2em]
    \item Definiciones clave se destacan en cajas \texttt{importante}.
    \item Teoremas o resultados importantes se derivan paso a paso.
    \item Ejemplos y ejercicios se incluyen para reforzar la comprensión.
\end{itemize}

\begin{importante}
Tip: Usa estas notas para repasar las deducciones analíticas del problema relativo, que permite pasar de 12 variables (6 por cuerpo) a solo 6 variables relativas más el movimiento uniforme del centro de masa.
\end{importante}

